\section{The comparison operator}\label{sec:comparison}

The factorisation of \cref{thm:factorisation} reduces the height of $Z(1)/Z(\xi_\iota)$ to that of the ratios $Z_j(1)/Z_j(\xi^j_\iota)$ of single lifts. To bound the latter, the
twisted lift $\Theta_{j,\vartheta}$ is compared with a twist-free map through an operator $U_\vartheta$ on the forms on $V_j$ (\cref{thm:comparison-operator}):
$\langle\Theta_\vartheta y',U_\vartheta\Theta_\vartheta y\rangle$ is a ratio of Hecke values times a fixed positive form, $U_\vartheta$ keeps integrality away from $Q$, takes lifts to forms orthogonal to the
false lifts, and $U_1$ is short. The operator is $U_\vartheta=n_{\rm Aut}h_k^2\,\mathrm{glue}_\vartheta\,\mathrm{glue}_\vartheta^*$, where $\mathrm{glue}_\vartheta$ builds forms on $V_j=V''\oplus V^\circ$ from tuples of forms on a
smaller definite space $V^\circ$ across a change of lattice, using \emph{first factors} on $V''$ obtained from three applications of the exchange identity, after which only a
number is left of the twist.

The section is organised as follows. The three exchanges turn a twisted theta series on $\U(W_j)$ into the twist-free functions $\mathsf F_i$, the first two by producing a fixed form, the carrier, on a space of
signature $(2m,m)$ (\cref{ss:three-exchanges}). The first factors and the gluing are defined and their integrality and sizes bounded in \cref{ss:gluing}. That the resulting map $\Theta^\circ$ is injective is a spanning
statement for products of theta series, proved by an isolation argument at the peaks of a form (\cref{ss:products-span}). The comparison operator is assembled from a tight frame in \cref{ss:comparison-operator}.

\subsection{The three exchanges}\label{ss:three-exchanges}
Fix $1\le j\le m$. Let $S$ be the space of signature $(2m,m)$, $I_n$ the standard definite space of rank $n$, $W^\flat=I_{6m+1}$, $V''$ the definite space on the indices of $W_j$, $(I_j)^-$,
$I_{3m-j}$ and $W^\flat$, and $V^\circ$ the definite space on the $r$ indices of $W_j$ and $t-9m-1-r$ further ones. The \emph{carrier} $\varphi^\circ$ is the theta lift of the constant function $1$ from the
unitary group of the negative definite $I_m^-$ to $\U(S)$, with the polynomial $\prod_\tau\det(y^-_\tau)^{2m}$; the \emph{carrier coefficients} $\mathsf a(h)c(\det h)$ are class sums over $\U(I_m)$ of theta kernels with
characters $c$, $1$. A \emph{choice} $i$ is a $3m\times\operatorname{rk}V''$ matrix $\Lambda$ modulo $\varpi^2$ with a class of $\U(W^\flat)$ at level $\varpi^2$ and a selection of rows; the function $\mathsf F_i$ on $\U(W_j)$ is the
integral over $[\U(S)]$ of the carrier against the theta kernel of $(S,W_j\oplus\mathrm{padding}\oplus W^\flat)$ with the trace character of $\Lambda$ as finite datum and a padded product of minors as
polynomial. The \emph{exchange pairs} of $j$ are three critical pairs: $(\mu^0_j\upsilon^n,z_*)$ with $n=j+(t-11m-1)/2$ and the main point $z_*$, and $(\upsilon^{-z_\pm}\chi^{\pm j},z_\pm)$ for the two signs. The
\emph{Hecke value product} $\Gamma_j(\vartheta)$ is the product over the three pairs of $L_*$ at the inverse characters twisted by $\vartheta$ or $\vartheta^{\pm1}$, for the chosen places and $Q$.
\lean{New.Rows.span_piM, New.Rows.stackedDet, New.Rows.unitModulus, New.Rows.carrierForm, New.Rows.carrierCoeff, New.Data.untwistedModulus, New.Data.ideal_untwistedModulus_eq_span, New.Data.CarrierIndex, New.Data.FlatIndex, New.Data.SecondFactorIndex, New.Data.SecondFactorClasses, New.Data.rowIndexOfNumber, New.Data.RowChoice, New.Data.shortRows, New.Data.twistedChiMPow, New.Data.twistedChiMPow_hasWeight, New.Data.PadIndex, New.Data.FirstFactorIndex, New.Data.padSigns, New.Data.negIndex, New.Data.posRowIndex, New.Data.untwistedDescentPair, New.Data.DetIndex, New.Data.paddedMinorProd, New.Data.firstFactorCharExp, New.Data.carrierChar, New.Data.carrierChar_hasWeight, New.Data.firstFactorChar, New.Data.ChoiceIndex, New.Data.choiceFn, New.Data.card_detIndex_add_firstFactorCharExp, New.Data.nStar, New.Data.two_mul_r_le_b, New.Data.mul_nStar_add_t_le_halfGap, New.Data.nStar_pos, New.staircaseWeight_of_eq_one, New.staircaseWeight_of_eq_neg_one, New.Data.orderW, New.Data.conjHolomorphicWeight, New.Data.kernelWeight, New.Data.kernelWeight_of_eq_one, New.Data.kernelWeight_of_eq_neg_one, New.Data.six_mul_r_le_n, New.Data.hundred_mul_m_add_nine_lt_k, New.Data.t_eq_five_mul_r, New.Data.card_firstFactorIndex, New.Data.two_mul_card_detIndex_add_le, New.Data.conjHolomorphicWeight_eq_add_kernelWeight, New.Data.conjHolomorphicWeight_of_eq_one, New.Data.conjHolomorphicWeight_of_eq_neg_one, New.Data.mul_nStar_le_sum_colsMinus, New.Data.abs_conjHolomorphicWeight_le, New.Data.isKDominant_kernelWeight, New.Data.isKDominant_conjHolomorphicWeight, New.Data.hasHolGap_kernelWeight, New.Data.hasHolGap_conjHolomorphicWeight, New.Data.isPosFirstOrder_orderW, New.Data.paramW_lt_of_orderW_lt, New.Data.fact_piM_pow_ne_zero, New.Data.places_twistModulus, New.Data.isTameFor_of_hasWeight, New.Data.isTameFor_twistedChiMPow, New.toPoly_coeff_nablaBasis, New.schur_mulVec_coeff_nablaBasis, New.Rows.exists_int_weylRep_invariant_fockPoly_eq_smul_monomial}
\lean{New.Data.signPoint, New.Data.mainPairPoint, New.Data.k_add_six_le_t, New.Data.signPoint_mem_Icc, New.Data.mainPairPoint_mem_Icc, New.Data.k_add_t_add_one_le_liftWeight, New.Data.liftWeight_le_four_mul_b, New.Data.exponent_add_mainPairPoint_le_liftWeight, New.Data.chiMLocal_zpow_eq_one_of_not_isChosen, New.Data.exchangePair, New.Data.heckeValueProd, New.Rows.exists_mul_stdClassSum_modKernelStdSpace_eq_one, New.Rows.exists_stdClassSum_modKernelStdSpace_neg_eq_one, New.Rows.exists_mul_stdClassSum_eq_carrierForm, New.Rows.carrierForm_mem_rapidOfLevel_and_isAdelicHol, New.Rows.fockInner_realPairWeil_variableMinor_pow_eq_self, New.Rows.carrierForm_ne_zero, New.Data.enlargedSpace, New.Data.enlargedPoint, New.Data.signOrder, New.Data.signs_eq_blockSigns_signOrder, New.Data.invariantTraceChar, New.Data.invariantTraceChar_mul, New.Data.twistedChiMod, New.Data.thirdExchangeKernel, New.Data.thirdExchangeLeftSide, New.Data.firstFactorChar_mul_upsilonU1_pow, New.Data.isTameFor_twistedChiMPow_mul_carrierChar, New.Data.matchedPoly_constantRows, New.Data.integral_carrierForm_mul_eq_card_mul_choiceFn, New.Data.exists_mul_thirdExchangeLeftSide_eq_choiceFn, New.Data.productDatum, New.Data.stdClassSum_thirdExchangeKernel_productDatum, New.Data.stdClassSum_thirdExchangeKernel_eq_productDatum, New.Data.thirdExchangeLeftSide_eq_kernelClassSum, New.Data.goodForms_mem_rapidOfLevel, New.Data.polyVector_letterPoly_mem_holomorphicData, New.Data.exists_choiceFn_eq_mul_kernelClassSum, New.Data.mul_choiceFn_mem_rapidOfLevel, New.Data.isAdelicHol_choiceFn}
\lean{New.Rows.FixesColsOff, New.Rows.rowStabilizer, New.Rows.columnDatum, New.Rows.columnDatum_mul_ite_eq_ite_mem_rowStabilizer, New.Rows.isRapid_modKernel, New.Rows.hasLeftInv_of_isUnit_det_toRows₁, New.Rows.areMatched_det_mul, New.Rows.integrable_and_mul_stdClassSum_eq_integral_mul, New.Rows.mul_stdClassSum_eq_integral_mul_modKernel, New.Rows.integral_mul_modKernel_sum_translates, New.Rows.stdClassSum_tower, New.Rows.signedRowSubst, New.Rows.realPairWeil_ratEmbed_inv_polyVector, New.Rows.modKernelStdSpace_stdEmbFin_mul, New.Rows.fixesLeft_modKernelStdSpace, New.Rows.stdClassSum_rowStabilizer_eq}

Here the numbers of the data are: the rank $r$ of $W_j$, $k=r+1-m$, $t=5r$, the gap $\Delta=2r^8$, $\Delta'=\Delta/2-(t+1)/2$, $r_+=k+j-1$ the number of positive signs of $W_j$; the points are $z_+=m$, $z_-=m/2$ and
$z_*=7(m/2)+k/2$. The \emph{characters} are $\chi'=\chi^j\vartheta$, $\nu''=\mu^0_j\vartheta\,\upsilon^{j+\lfloor(t+k-1)/2\rfloor}$, the character of the $j$-th lift times a power of $\upsilon$, and the \emph{carrier character}
$\nu^\flat\upsilon^m=\nu''\upsilon^m/\chi'$, which is free of the twist and has trivial tame part. The \emph{padded polynomial} is $P=\prod_\tau\det(y^-_\tau)^{e'(\tau)}\prod_{h<r_+}\det(y_{\tau,h})^{k^+(\tau,h)}$ on $(S,W_j\oplus I_j^-\oplus I_{3m-j})$, with
$y^-_\tau$ the $m\times m$ block in the $m$ chosen negative rows and $y_{\tau,h}$ the $2m\times2m$ block in the $2m$ rows starting at the basis vector of number $h$ of $W_j$, and $k^+$ the column counts of the diagram
$\mu^+$ of the $j$-th lift; on $W^\flat$ the polynomial is the product of determinants pairing the chosen rows of a short basis of the class with the positive variables. The index set $\mathrm{Det}_\tau$ has one
element for each column of $\mu^+$ at $\tau$: $k^+(\tau,h)$ for each height $h<r_+$ and $\Delta'$ of height $r_+$.

\begin{lemma}[Numerics of the exchanges]\label{lem:exchange-numerics}
\begin{enumerate}[label=\textup{(\alph*)}]
\item $\#\mathrm{Det}_\tau+j+\lfloor(t+k-1)/2\rfloor=\kappa_j(\tau)$; $\#V''=r+9m+1$; $6r\le n_m$, $100m+9<k$, $k+6\le t$, $2r\le b$, and $2(\#\mathrm{Det}_\tau+\#V''+m)\le5b$.
\item $z_\pm,z_*\in[1,t]$; $k+t+1\le\kappa_j\le4b$; and $n+z_*\le\kappa_j$ for $n=j+(t-11m-1)/2$. Hence the three exchange pairs are critical pairs, ramified at the chosen places.
\item With $n_*=2ft^3$, $300n_*+t\le r^8$.
\item For $c\ne0$ and a twist $\vartheta$, a character of weight $\gamma$ with tame part $\psi\cdot\mathrm{tame}(\vartheta)$, $\psi$ nontrivial exactly at the chosen places, is tame for $\varpi^cQ$; in particular so is $\chi'$, and so is $\chi'\nu^\flat\upsilon^m$, of
weight $m-\#\mathrm{Det}_\tau$, with the residue character of $\chi'$.
\end{enumerate}
\end{lemma}

\begin{proof}
(a) For $h<r_+$ the number of columns of height $h$ is the successive difference of the shifted target parameter minus one, so the sum telescopes, and $k$ is even and positive; the other statements compare the numbers of
the data. (b) The first two are comparisons of the data, the weights $\kappa_j=\kappa_0+j\Delta$ growing by the gap; the third is the same comparison with the explicit values. The weights of the pairs are $-(\kappa_j-n)$ and $-z_\pm$,
between their points and $4b$, and their tame parts $\omega^j$, $\omega^{\pm j}$ are nontrivial exactly at the chosen places, $\omega$ having order greater than $8m$ there. (c) is a comparison of powers of $r$. (d) The tame part of $\psi$ is
nontrivial exactly at the chosen places and that of $\vartheta$ exactly at $Q$; the two sets are disjoint and their union is the set of places of $\varpi^cQ$. The weight of $\chi'\nu^\flat\upsilon^m$ is that of $\nu''\upsilon^m$ minus $0$, which is
$m+j+\lfloor(t+k-1)/2\rfloor-\kappa_j=m-\#\mathrm{Det}_\tau$ by (a).
\end{proof}

For the one exchange, let $W=I_b\oplus W'$, $\mathfrak m\ne(1)$, $\chi$ a character of $\OO_E/\mathfrak m$, $e\colon r\hookrightarrow N$ and $w$ a family of rows modulo a quotient of the residue ring. The \emph{row stabiliser} $H$ is the group of invertible $A$
whose columns outside $e$ are those of $1$ and with $w\,\phi(A)=w$; the \emph{column datum} of $\sigma$ is $\chi$ of the determinant of the rows $e$ of $\sigma$ if $w\phi(\sigma)=w\phi(\iota_e)$, and $0$ otherwise.

\begin{lemma}[One exchange and the row stabiliser]\label{lem:one-exchange}
\begin{enumerate}[label=\textup{(\alph*)}]
\item If $F$ is left invariant under $H$, then the column datum of the columns $e$ of $z$, times $F(y)$ when $z$ has the columns of $1$ off $e$ and $0$ otherwise, equals $\chi(\det z)F(z^{-1}y)$ for $z\in H$ and $0$ for the other $z$.
\item If $\Psi'$ is invariant under right multiplication by invertible matrices, then $\Phi(x'\ x'')=\chi(\det x'')\Psi'(x''^{-1}x')$ and $\Psi\binom{y'}{y''}=\chi(\det y')\Psi'(y'')$ are matched for $\chi^{-1}$.
\item A theta kernel whose finite datum lives on left invertible matrices decays rapidly in the adelic point.
\item (One exchange.) If $c$ times the class sum of a kernel family $K\Phi$ is the integral over $[\U(V')]$ of the kernel of $(V',W)$ with datum $\Psi$ and polynomial $P$ on $W'$ for all matched $\Phi$, $\Psi$, then for $\Psi'$ as in (b) the
product of the kernel of $(V',I_b)$ at $h'$ with datum $\chi\circ\det$ and the kernel of $(V',W')$ at $g$ with datum $\Psi'$ is integrable, with integral $c$ times the class sum of $K\Phi$ for $\Phi$ as in (b); the same summed over the classes of
$\U(I_b)$ against coefficients.
\item (Tower.) If $f_0$ is left invariant under the rational points and right invariant under the integral points reducing into $H$, and $f$ is the sum of the right translates of $f_0$ by lifts of coset representatives of $H$, then
the class sum of $f_0$ at the level of $H$ is the class sum of $f$ at full level.
\item Rational points of $\U(I_b)$ act on a kernel of $(W,I_b)$ by a substitution of the polynomial times a power of the determinant; if the Fock vector of $P$ is fixed by the real group of $I_b$, the kernel is left invariant
under the rational points. Under these hypotheses the class sum at full level of the kernel with datum $\chi(\det x'')F(x''^{-1}x')$ equals the class sum at the level of $H$ of the kernel with the restricted datum, the column datum
of $\sigma$ times $F(x')$.
\end{enumerate}
\end{lemma}

\begin{proof}
(a) On a matrix with the columns of $1$ off $e$, the column datum can be read on the whole matrix: for $z=1+u\iota_e^{\mathsf T}$ the columns $e$ are $\iota_e+u$ and their minor on the rows $e$ is $\det z$; both conditions on $w$ say
$w\phi(u)=0$. For $z\in H$ the invariance of $F$ gives $F(z^{-1}y)=F(y)$. (b) The block determinant functions of \cref{prop:matched}(a) for the function $f=\Psi'$ give the match $\chi(\det y')\Psi'(-y''y'^{-1})$, and the invariance of $\Psi'$ removes
the factor $-y'^{-1}$. (c) The datum vanishes on the matrices that are not left invertible modulo $\mathfrak m$, so in the theta series only vectors of full rank occur; high in a cusp of the first space these
are long, and the kernel is at most a constant times a Fock moment times $\mathrm{ht}(h)^B\mathrm{ht}(g)^{-M}$, for every $M$. (d) The function matched with
$\chi(\det x'')\Psi'(x''^{-1}x')$ is the product $\chi(\det y')\Psi'(y'')$ by (b), so its kernel factors along $I_b\oplus W'$; the integral converges because $y'$ is invertible on the support, by (c). The class sum is finite, so it commutes
with the integral. (e) The number of cosets cancels the index of the level. (f) The kernel is invariant under the rational point $\delta^{-1}$, which cancels the finite translation and leaves the real part of $\delta^{-1}$, acting on the
Fock vector by the substitution $x_\tau\mapsto(\delta^{-1})^{\mathsf T}x_\tau$ in the positive columns and $x_\tau\mapsto\delta x_\tau$ in the negative ones, times $\prod_\tau\det\delta^{q-\gamma_\tau}$. For the last claim, the datum $\Phi_0$ is $\Phi$ cut down to
$x''\in H$, by (a); the twisted translates of $\Phi_0$ by coset representatives of $H$ add up to $\Phi$, an invertible $x''$ lying in exactly one coset, and their kernels are right translates of the kernel of $\Phi_0$, the twist cancelling the
factor that the translation adds to $\chi(\det x'')$. (e) concludes.
\end{proof}

\begin{lemma}[The carrier]\label{lem:carrier}
Let $m=2e>0$.
\begin{enumerate}[label=\textup{(\alph*)}]
\item (First exchange.) For exponents $k_\tau$ there is $C$ such that each $\alpha$ of weight $-k$ tame for $(\mathfrak m,\chi)$ has $\gamma$ with $\abs\gamma^2=C/L_*(\alpha^{-1}\upsilon^e,e)$ and $\gamma\sum_h\theta_{\alpha,\alpha}[\chi\circ\det,\prod_\tau\det(y_\tau)^{k_\tau}](g,h)=1$ on
$\U(I_m)$, the class sum at the stabiliser of the standard lattice.
\item (Conjugation.) There is $n$ of modulus one such that if $\gamma$ times the class sum on $I_m$ of the kernel of $(I_m,I_m)$ with characters $c^{-1}$, $c^{-1}$ and datum $\chi^{-1}\circ\det$ is $1$, then so is the conjugate of $\gamma$ times
$n/\chi(\eta)^m$ times the class sum of the kernel of $(I_m^-,I_m)$ with characters $c\upsilon^m$ and datum $\chi\circ\det$, $\eta$ the reduction of an element of norm $-1$.
\item The vector of $\prod_\tau\det(y^-_\tau)^{2m}$ in the Fock space of $(I_m^-,S)$ is fixed in the matrix coefficient: $\langle v,\Omega(u,1)v\rangle=\langle v,v\rangle$ for $u$ in the compact real group of $I_m^-$; hence the carrier is not zero, for every $m>0$.
\end{enumerate}
\end{lemma}

\begin{proof}
(a) Apply \cref{cor:explicit-exchange} with $W=I_m$ and $V'=0$: the data $\chi\circ\det$ and $1$ are matched (the block determinant functions for the constant $1$), the integral over the trivial group is $1$, the polynomial of the
empty matrix is $1$, and with no rows to add the stacked determinant product is the product of powers of determinants; up to the factor $\alpha(\det x)^{-1}$ the kernel with the zero space has the standard characters.
(b) Conjugating the kernel replaces the space by its opposite and the datum by its conjugate; it scales $\chi\circ\det$ by $(\chi(\eta)^m)^{-1}$, and kernels are linear in the datum. (c) Every variable of a minor in the negative rows has sign
$1$, each $u$ scales the polynomial by the $2m$-th powers of its determinants, and these are products of powers of determinants of unitary matrices, of modulus one; the coefficient is the constant $\langle v,v\rangle$. If the lift
vanished, the real convolution of the coefficient would vanish (Rallis' injectivity argument, \cref{thm:B-rallis}(c), the finite part being the vacuum, which is sharp), but the integral of a nonzero constant over the compact group
is not zero.
\end{proof}

\begin{proposition}[The carrier and the three exchanges]\label{prop:three-exchanges}
\begin{enumerate}[label=\textup{(\alph*)}]
\item Let $m$ be even and let a tame character of weight $0$ exist. Then the carrier is continuous, left invariant under the rational points and of every principal level; it is rapidly decreasing; and
it is holomorphic for an ordering of the basis with the positive vectors first.
\item (First two exchanges.) Let $m=2e$. There is $C$ such that for each character $c$ of weight $0$ tame for $(\mathfrak m,\chi)$ some $\gamma$ with
\[
  \abs\gamma^2=\frac{C}{L_*(c^{-1}\upsilon^m,m)\,L_*(c\upsilon^e,e)}
\]
satisfies
$\gamma\sum_h\mathsf a(h)c(\det h)\,\theta_{\upsilon^m,c}[\chi\circ\det,1](x,h)=\varphi^\circ(x)$.
\item (Third exchange.) There are $C$ and numbers $c(\vartheta)$ with $\abs{c(\vartheta)}^2=C/\Gamma_j(\vartheta)$ such that $c(\vartheta)$ times the left side of the third exchange, a theta series on $\U(W_j)$ (a class sum of kernels
of $(W_j,V'')$ with the vectors $\Upsilon_{\vartheta,i}(h)$), is $\mathsf F_i$ for every choice $i$.
\item (One exchange.) For $W=I_b\oplus W'$ and a datum invariant under a row stabiliser, $c$ times the class sum of the kernel with datum $\chi(\det x'')\Psi'(x''^{-1}x')$ is the integral of the kernel of
$(V',I_b)$ with datum $\chi\circ\det$ times the kernel of $(V',W')$ with datum $\Psi'$; a class sum at full level equals a class sum at the level of the row stabiliser.
\end{enumerate}
\end{proposition}

\begin{proof}
(b) Chain two exchanges by \cref{lem:one-exchange}(d). The first is \cref{lem:carrier}(a) for $c^{-1}$, conjugated by \cref{lem:carrier}(b) so that it lives on $I_m^-$: the class sum on $I_m$ of kernels of $(I_m^-,I_m)$ is constant on $\U(I_m^-)$.
The second is \cref{cor:explicit-exchange} with $W=I_m\oplus S$, $V'=I_m^-$ and $V=I_{3m}$, the datum $\Psi'=1$, summed over the classes of $I_m$: a multiple of a double class sum of kernels of $(W,I_{3m})$ is the integral of the kernel of
$(I_m^-,S)$ with datum $1$, which is the carrier. The character of the second exchange is $(c\upsilon^m)^{-1}\upsilon^{2m}=c^{-1}\upsilon^m$. The kernel on $I_m\oplus S$ splits into the summand of the carrier coefficient and the kernel on $S$, and both sides
are the claim times $c(\det x)$. The constants $n$ and $\chi(\eta)$ have modulus one, so the squared moduli of the two constants multiply to the stated value.

(a) By (b), applied to the tame character of weight $0$, the carrier is a finite combination of theta kernels with polynomial $1$; level, rapid decay and holomorphy are read off from these kernels. Level one: on the stabiliser of the
standard lattice the factor $c(\det k)^{-1}$ cancels against $\chi(\det k)$. Rapid decay: $\chi(\det y)\ne0$ makes $y$ invertible, so \cref{lem:one-exchange}(c) applies. A theta kernel is continuous and invariant under the rational points.
For holomorphy, order the basis with the positive indices first; the polynomial $1$ is a product of zeroth powers of minors, so its Fock vector is holomorphic, and so is every kernel with it.

(c) \emph{The exchange.} Apply \cref{cor:explicit-exchange} to $W=I_{3m}\oplus W'$, where $W'=W_j\oplus I_j^-\oplus I_{3m-j}\oplus W^\flat$ has the index set of $V''$, to $V''$ and to the space $S$ (as $V'$), and sum it over the classes of $I_{3m}$ by
\cref{lem:one-exchange}(d); by (b) the square kernel of $S$ becomes the carrier. The basis of $W$ is listed with the positive vectors first, the $m$ negative ones forming the second block, so the signs are those of a block
decomposition. For the constant rows of a choice $i$, the matched polynomial of \cref{cor:explicit-exchange} involves only the rows of $W'$: it is the padded product $P$ times the determinants on $W^\flat$ formed with the rows of $i$, by
comparing the exponent of the negative minor, the positive minors and the determinants on $W^\flat$ entry by entry. \emph{The datum.} The left side carries the datum $(x'\ x'')\mapsto\chi'(\det x'')\psi^{\rm inv}_\Lambda(x''^{-1}x')$, where
$\psi^{\rm inv}_\Lambda$ is the sum over $\GL_{3m}(\OO_E/\varpi^2)$ of the translates of the trace character of $\Lambda$; it is invariant under right multiplication, as \cref{lem:one-exchange}(b) requires. The carrier has full level, so every translate
gives the same integral, and the integral of the carrier against the kernel with $\psi^{\rm inv}_\Lambda$ is $\#\GL_{3m}(\OO_E/\varpi^2)$ times $\mathsf F_i$. \emph{The characters.} By \cref{lem:exchange-numerics}(d) and the identities
\[
  \frac{(\nu^\flat\upsilon^m)^{-1}}{\chi'}\upsilon^{3m+z_*}=\frac{\vartheta^{-1}}{\mu^0_j\upsilon^n},\qquad\chi'^{-1}\upsilon^m=\frac{\vartheta^{-1}}{\upsilon^{-m}\chi^j},\qquad\chi'\upsilon^{m/2}=\frac{\vartheta}{\upsilon^{-m/2}\chi^{-j}},
\]
obtained by comparing exponents, the characters of the three exchanges are the inverses of those of the three pairs twisted by $\vartheta$ or $\vartheta^{\pm1}$, so the squared moduli of the three constants multiply to a constant over
$\Gamma_j(\vartheta)$. \emph{The theta series.} If the rows of $\Lambda$ are combinations of rows $w$ modulo $\varpi^2$, $\psi^{\rm inv}_\Lambda$ is invariant under the matrices fixing $w$, and the real group of $V''$ fixes the polynomial of the kernel;
by \cref{lem:one-exchange}(f) the class sum at full level equals the class sum at the level cut out by $w$ of the kernel with the product datum, the column datum of $w$ times $\psi^{\rm inv}_\Lambda(x')$. Expanding the series $\Upsilon$ in
coordinates, each term is one kernel on the index set of $W_j\oplus I_j^-\oplus W^{\rm add}$, with $W^{\rm add}=I_{3m-j}\oplus W^\flat\oplus I_{3m}$, and relabelling matches the left side; so the left side is the class sum over $\U(V'')$ of kernels of
$(W_j,V'')$ with the characters of the descent pair and $\nu''$, the column datum, and the polynomials of the vectors $\Upsilon_{\vartheta,i}(h)$.

(d) is \cref{lem:one-exchange}(d), (f).
\end{proof}

\subsection{First factors and gluing}\label{ss:gluing}
The \emph{series} $\Upsilon_{\vartheta,i}$ on the finite points of $\U(V'')$ is built from theta kernels of definite spaces and the datum $\Psi^{\rm add}$ on $W^{\rm add}=I_{3m-j}\oplus W^\flat\oplus I_{3m}$; the \emph{first factor}
$\varphi''_{\vartheta,i}$ is the complex conjugate of $c''\Upsilon_{\vartheta,i}$ for a constant $c''\ne0$, and $\Xi_{\vartheta,i}(g)=\langle\varphi''_{\vartheta,i},\theta''_{\vartheta,\Lambda}(g,\cdot)\rangle$. A form is \emph{integral and small} if its values are integral away from the places of
$Q$ and, if $Q$ is empty, their conjugates have squared norm at most $\exp(fmb\log^2b)$. The \emph{gluing} $\mathrm{glue}_\vartheta$ takes a tuple $\zeta$ of forms on $V^\circ$ (one copy for each choice, given by values at
the classes of level $\varpi^2$) to the form on $V_j=V''\oplus V^\circ$ supported, choice by choice, on the classes of $(h,h^\circ)\mathsf g$ for a lattice $\mathsf g$ glued for $\Lambda$, with values the Cartan products of
$\varphi''_{\vartheta,i}(h)$ and $\zeta(i,c)$. Write $\Theta^\circ=c'_1\mathrm{glue}_1^*\Theta_{j,1}$ on good forms.
\lean{New.HasWeightNormSqLe, New.HasWeightNormSqLe.nonneg, New.HasWeightNormSqLe.one, New.HasWeightNormSqLe.mul, New.HasWeightNormSqLe.prod, New.HasWeightNormSqLeIf, New.hasWeightNormSqLeIf_iff, New.HasWeightNormSqLeIf.prod, New.eval_weylKernel_single_eq_det, New.exists_toPoly_single_eq, New.exists_toPoly_eq_det_and_weylNormSq_le, New.Data.RowChoice.classRep, New.Data.RowChoice.rows, New.Data.ChoiceIndex.matrix, New.Data.ChoiceIndex.rowChoice, New.Data.AddedIndex, New.Data.regroupEquiv, New.Data.posEquiv, New.Data.constantRows, New.Data.stackedDetProdCoords, New.Data.forgetTwist, New.Data.addedDatum, New.Data.leftSideSeries, New.Data.card_detIndex_eq_sum, New.Data.card_detIndex_le, New.Data.card_detIndex, New.Data.prod_sum_norm_sq_le_M_pow, New.Data.IsSmallSpecialization, New.Data.hasWeightNormSqLeIf_stackedDet, New.Data.hasWeightNormSqLeIf_stackedDetProd, New.Data.comp_eval_stackedDetProdCoords_eq_of_toPoly_eq, New.Data.toPoly_comp_eval_stackedDetProdCoords, New.Data.weylNormSq_comp_eval_stackedDetProdCoords_le, New.Data.weylNormSq_stackedDetProdCoords_le, New.Data.toPoly_stackedDetProdCoords, New.Data.aeval_smul_X_stackedDetProdCoords, New.Data.inFirstCols_stackedDetProd_mem_range, New.Data.IsHomogeneousOverEntries, New.Data.exists_isHomogeneousOverEntries_smul_stackedDetProdCoords}
\lean{New.Data.renameRowsCols_stackedDetProd_eq_sum, New.Data.normQ_mul_M_pow_mem_modulus, New.Data.constantRows_mul_transpose_eq_mul_fromBlocks, New.Data.firstFactorSize_sq_le_exp, New.Data.modKernelBilin_stackedDetProdCoords_equivariant, New.Data.norm_shortRows_mul_transpose_le_one, New.Data.weylGauge_modKernelBilin_stackedDetProdCoords_le_one, New.Data.addedDatum_mem_integralAway, New.Data.norm_map_addedDatum_le, New.Data.card_le_gap, New.Data.sqrt_weylNormSq_smul_eval_stackedDetProdCoords_le, New.Data.sqrt_weylNormSq_comp_le_of_eval_le, New.Data.weylGauge_stdClassSum_stdClassSum_le_one, New.Data.weylNormSq_stdClassSum_stdClassSum_le, New.Data.IsIntegralAndSmall, New.Data.norm_mul_modKernelStdSpace_twistedChiMPow_le_one, New.Data.norm_mul_modKernelStdSpace_one_le_one, New.Data.exists_star_leftSideSeries_eq_classSum, New.Data.exists_isIntegralAndSmall_leftSideSeries}
\lean{New.exists_det_fromCols_eq_mul_det, New.Rows.twoSidedSubst, New.Rows.inFirstCols, New.Rows.stackedDetProd, New.Rows.twoSidedSubst_stackedDetProd, New.Rows.firstColsCoord, New.Rows.twoSidedSubst_firstColsCoord}
\lean{New.LevelPlaces.chosenModulusSq, New.Rows.map_block_eq_blockGL, New.Rows.levelModulus_ideal_le, New.Rows.places_levelModulus_false, New.Rows.reduceToChosenModulusSq, New.Rows.liftMod_mem_supportedOn, New.Rows.atPlace_twistPlace_eq_one, New.Rows.DividedSplitting, New.DescentPair.genIdele, New.DescentPair.genIntAdele, New.Rows.genIdele_apply_of_notMem, New.Rows.reduction_levelModulus_eq_zero_iff, New.Rows.DividedSplitting.genResidue, New.Rows.DividedSplitting.genResidue_smul_eq_inclMatrix_iff, New.Rows.kernelVector, New.Rows.rowFixingLevel, New.Rows.kernelClassSum, New.Rows.DividedLattice, New.Rows.DividedLattice.reducedInverse, New.Rows.DividedLattice.weilScalar, New.Rows.DividedLattice.reduceAt_conj_eq_one, New.Rows.DividedLattice.exists_reducedInverse_eq, New.Rows.DividedSplitting.stackOverGenScalar, New.Rows.rowFixingLevel_le_levelOf, New.Rows.DividedSplitting.congruenceGroup, New.Rows.DividedSplitting.mem_congruenceGroup_iff, New.Rows.DividedSplitting.reduceAt_liftMod_congruenceGroup, New.Rows.DividedSplitting.IsBlockNormalised, New.Rows.DividedSplitting.IsGlued, New.Rows.DividedSplitting.IsGlued.indicator_eq_sum, New.Rows.DividedSplitting.productForm, New.Rows.DividedSplitting.exists_congruenceGroup_inv_eq_one_sub, New.Rows.DividedSplitting.glueMatrix, New.Rows.DividedSplitting.isBlockNormalised_iff, New.Rows.DividedSplitting.reducedInverse_mul_eq_iff, New.Rows.DividedSplitting.mul_congruenceGroup_eq_smul_iff, New.Rows.DividedSplitting.isBlockNormalised_mul_iff, New.Rows.DividedSplitting.exists_divisor_assoc_stackOverGenScalar, New.Rows.DividedSplitting.exists_genResidue_mul_stackOverGenScalar_eq_zero, New.Rows.DividedSplitting.exists_mul_glueMatrix_eq, New.Rows.DividedSplitting.exists_isGlued}
\lean{New.Rows.kernelClassSum_smul, New.Rows.letterPoly_schur_inclMatrix_mulVec, New.Rows.fockPoly_cartanProduct, New.Rows.DividedSplitting.idem_mul_mem_of_inv_mul_mem, New.Rows.DividedSplitting.liftFn_datum_point_inv_mul, New.Rows.DividedSplitting.normalizedDescent_classForm_mul, New.Rows.DividedSplitting.prod_twist_det_eq_inv_mul_charOfTamePart, New.Rows.DividedSplitting.normalizedDescent_productForm, New.Rows.DividedSplitting.normalizedDescent_stdClassSum_productForm, New.Rows.DividedSplitting.denominatorExponent, New.Rows.DividedSplitting.norm_schur_diagonal_piecewise_inv_le, New.Rows.DividedSplitting.norm_schur_inv_cols_le_weylGauge, New.Rows.DividedSplitting.weylGauge_cartanProduct_le}

The \emph{constant rows} of a choice $i$ form, for each column of $\mu^+$ at $\tau$, a matrix $\mathsf c$ with $2m$ rows and with columns the $r_+$ positive indices of $W_j$ and the indices of $W^{\rm add}$: for a column of height $h<r_+$, row $a$ is the
basis vector of $W_j$ of number $h+a$ if $h+a<r_+$ and the basis vector number $a$ of $I_{3m-j}$ otherwise; for the $l$-th column of height $r_+$ the rows are the chosen rows of $\tau(Y)$, for a short basis $Y$ of the class of $i$,
placed in the columns of $W^\flat$. The polynomial $P_V$ is the product of the determinants $\det(\mathsf c;Xz)$ of the constant rows stacked over the rows of variables indexed by $V''$, with $X$ the variables in the positive columns of
$W_j$ and $z$ those in $W^{\rm add}$. The datum $\Psi^{\rm add}$ on matrices $z$ modulo $\varpi^2Q$ is $0$ unless the columns on $I_{3m-j}$ and $W^\flat$ are those of $1$, and then is $\psi^{\rm inv}_\Lambda(z')$ for the columns $z'$ on $I_{3m}$. A
polynomial $F$ \emph{has weight $\kappa$ and squared norm at most $b$} if it is the polynomial $\sum_i(G\xi)_ib_i$ of a vector $\xi$ of the Weyl module of $\kappa$ (with $G$ the Gram matrix and $b_i$ the products of minors of the standard
fillings) with $\xi^*G\xi\le b$.

\begin{lemma}[Weights and norms of determinants]\label{lem:det-weights}
\begin{enumerate}[label=\textup{(\alph*)}]
\item Under products the weights add and the bounds on the squared norms multiply, and the same holds for bounds valid under a condition.
\item For the weight of one column of height $h$ at $\tau$, a function of $h$ columns that is the sum over the $h$-subsets of the minor on the subset times its value at the unit columns is the polynomial of a vector whose squared norm is
the sum of the squared moduli of these values; and $\det((0;X)\mid Z)$ is the polynomial of a vector of squared norm at most the product of the squared lengths of the columns of $Z$.
\item There is $Y$ with $\det(X\mid Z)=\gamma\det(YX)^{\mathsf T}$ for all $X$. The substitution $x_\tau\mapsto c_\tau^{\mathsf T}x_\tau g_\tau^{\mathsf T}$ sends the stacked determinant product for the matrices $M_{\tau,k}g_\tau^{\mathsf T}$ to that for $M_{\tau,k}$ with each factor
multiplied by $\det c_\tau\det g_\tau$; and if it fixes the first $p$ columns and sends $F'$ to $sF$, the coordinates of $F'$ in the first $p$ columns, combined by the Schur matrix, go to $s$ times those of $F$.
\end{enumerate}
\end{lemma}

\begin{proof}
(a) The product of two semi-invariants is a semi-invariant of the sum of the weights, and the Gram norm of the product vector is at most the product of the norms. (b) Evaluation at the unit columns on a subset has the minor on the
subset as polynomial, and the pairing of $\xi$ with this evaluation is the value; for the determinant, expand along the first $h$ columns: up to sign the coordinates are the complementary minors of $Z$, and the Cauchy--Binet formula with
Hadamard's inequality bounds the sum of their squared moduli. (c) Choose an invertible $(U_1\mid U_2)$ with $U_1^*Z=0$ and put $Y=U_1^*$; then $(U_1\mid U_2)^*(X\mid Z)$ is block triangular, and comparing determinants gives the first claim. For one factor
of the product with matrix $A$ on the right, the matrix on the left is $\mathrm{diag}(1,c)^{\mathsf T}Ag^{\mathsf T}$, which gives the second. For the third, read the polynomials in the first columns: the substitution by $(c,\mathrm{diag}(1,g))$ splits into the
substitution on the coefficients and the row substitution by $c^{\mathsf T}$ on the first columns, whose adjoint for the Fock pairing is the row substitution by $c$ (\cref{ss:B-schur}), acting on the basis by the Schur matrix.
\end{proof}

\begin{proposition}[First factors]\label{prop:first-factors}
\begin{enumerate}[label=\textup{(\alph*)}]
\item Let $P_V$ be the product of the determinants $\det(\mathsf c;Xz)$, one for each column of $\mu^+$. Then $P_V=P_{\zeta(z)}(X)$ for a polynomial map $\zeta$ to the Weyl module of $\mu^+$. The map $\zeta$ is
homogeneous of one degree $d$, one integer $H\le(d^d)^3$ clears its denominators, and its squared norm is bounded when the entries are small.
\item For one $c''\ne0$, independent of $i$, $\vartheta$ and $h$, every first factor is integral and small.
\item There are $C$ and $c'(\vartheta)$ with $\abs{c'(\vartheta)}^2=C/\Gamma_j(\vartheta)$ and $c'(\vartheta)\Xi_{\vartheta,i}=c''\mathsf F_i$ for all choices.
\end{enumerate}
\end{proposition}

\begin{proof}
(a) \emph{The map.} Let $\zeta$ be the vector of coordinates of $P_V$ in the first $r_+$ columns $X$: read $P_V$ as a polynomial in $X$ with coefficients in the polynomials in $z$, and take its coordinates in the basis of products of
minors of weight $\mu^+$. By \cref{lem:det-weights}(b) and (a), each factor $\det(\mathsf c;Xz)$, specialised in $z$, is the polynomial of a vector of the weight of one column of its height, so the product is the polynomial of a vector of
weight $\mu^+$: below $r_+$ each height occurs once for each column of $\mu^+$, and at $r_+$ it occurs $\Delta'$ times, which is the column count of the descent pair there. Coordinates commute with specialisation, so $P_V=P_{\zeta(z)}(X)$, and,
comparing coefficients at every point, $P_V$ is the polynomial of $\zeta$ identically. \emph{Degree and denominator.} Scaling the variables by $c$ multiplies the product of determinants by $c^{N\sum\kappa(h)}$, with $N=\#V''$, while the
Schur matrix of $c\cdot1$ is the scalar $c^{\sum h\kappa(h)}$; so each coordinate is homogeneous of degree $d=\sum(N-h)\kappa(h)$, and $N-h\le10m+1$ whenever $\kappa(h)\ne0$. The coefficients of $P_V$ in the first columns lie in every subring of $\C$
containing the conjugates of the entries of $Y$; the coordinates are obtained from them by the inverse Gram matrix, which $H$ clears, with $H$ the hook number of the complementary diagram $h\mapsto\kappa(N-h)$, which has at most $d$ boxes,
so $H\le(d^d)^3$. \emph{Norm.} If the conjugates of the entries of $Y$ and of the values of $z$ are at most $M^4$, each column of a determinant has squared length at most $(r+11m+1)M^8\le M^9$, at most $12m+1$ columns occur, and by
\cref{lem:det-weights}(b), (a) the squared norm of $\zeta(z)$ is at most $(M^9)^{(12m+1)N_{\rm det}}$ with $N_{\rm det}$ the total number of determinants, at most $\#\Phi(\lceil b\rceil+w_{\rm off}+\Delta+m\Delta)$ by \cref{lem:exchange-numerics}(a).

(b) \emph{The series as a class sum.} By the definition of the carrier coefficients, there is $c$ of modulus one, depending only on $j$, such that the conjugate of a multiple of $\Upsilon_{\vartheta,i}(h)$ is $n_{\rm Aut}^2$ times a double class sum
over $\U(I_{3m})$ and $\U(I_m)$ of conjugates of products of three theta kernels of definite spaces: $k_0$ of $(I_m,I_{3m})$ with characters $\chi'$ and $1$; the conjugate of $k_1$ of $(I_j,V'')$, whose datum is the conjugate, at an element of norm
$-1$, of the indicator of the inclusion matrix; and the kernel of $(W^{\rm add},V'')$ at the polynomial vector $H\zeta$, with datum $\Psi^{\rm add}$, both last kernels of modulus $\varpi^2Q$. \emph{Constants.} Each of the three kernels is made
integral, at points integral above $p\nmid\Nm Q$, by a constant (\cref{prop:mod-kernel}), the data being integral away from $\Nm Q$: the values of $\Psi^{\rm add}$ are $0$ or sums over $\GL_{3m}(\OO_E/\varpi^2)$ of values of an additive character, roots of
unity, and their conjugates are at most $\#\GL_{3m}(\OO_E/\varpi^2)$; $\Nm Q\cdot M^c$ lies in the ideal of $\varpi^cQ$. For $k_0$ and $k_1$ move the points into the integral points, the polynomials and the characters of weight $0$ not changing
the kernel. \emph{Integrality.} By weak approximation, rational points move the classes and $h$ into the integral points; by equivariance under rational points (the Schur matrix of $\gamma'^{-1}$ mixing the coordinates of $\zeta$ for a translated
short basis and the powers of the determinants cancelling against the change of the characters), the gauge at $h$ is that at an integral point, where the entries of $Y\gamma^{\mathsf T}$ are integral; each class sum takes one factor $n_{\rm Aut}$,
a common multiple of the orders of the stabilisers. \emph{Size.} For $Q$ empty, bound each term by duality: if the conjugates of the values of a vector of polynomials on a set are at most $B$ and a linear form $\theta$ has
$\abs{\sigma\theta(Q)}\le C$ for every homogeneous $Q$ small on the set, then the conjugate of $\theta$ composed with the vector is at most $C$; apply it to the third kernel with the bound of (a), multiply the three kernel bounds
$\exp(x_1)$ and count the classes, at most $\exp(x^{30})$ each. With $d\le(10m+1)N_{\rm det}\le M$ and $H\le d^{3d}\le M^{3d}$, the logarithm of the square of the product is at most $fmb\log^2b$ by the comparison of the sizes of the data.

(c) is \cref{prop:three-exchanges}(c), with $\Xi_{\vartheta,i}=c''$ times the left side of the third exchange: both are the class sum $\sum_h$ of kernels of $(W_j,V'')$ at the level $J_\Lambda$ with the datum given by the rows of $\Lambda$ stacked over
$\varpi\cdot1$ and polynomial $\varphi''_{\vartheta,i}(h)$. Put $c'(\vartheta)=c(\vartheta)$ of that proposition.
\end{proof}

The gluing works on a \emph{divided splitting} of $V=V''\oplus V^\circ$: the indices split as $N\sqcup m$, with $r$ indices of $W$ in $V''$ (the fixed embedding $e$), $r$ marked indices $z$ of $V^\circ$, and a copy of an index set $a$ in $N$; the diagrams $\mu^+$ and
$\mu^-$ have no columns of heights between the ranks of the summands and that of $V$. Let $\varpi$ be the idele generating the product of the chosen places. A \emph{divided lattice} is a point $\gamma$ of the finite unitary group, supported on the
chosen places, with $\varpi\gamma^{-1}$ integral (so $\varpi\OO^n\subseteq\gamma\OO^n\subseteq\OO^n$ there); its \emph{scalar} is $c_\gamma=\chi(\delta)/\sqrt{\abs\delta}^{\,r}$ with $\delta$ the determinant of $\gamma$ at the chosen places. The \emph{congruence group}
$\bar J^1$ consists of the invertible $r\times r$ matrices modulo $\varpi^2$ that are $1$ modulo $\varpi$. For a matrix $s$ modulo $\varpi^2$ with rows indexed by $a$, a matrix $x$ is \emph{block normalised} if its rows in $m$ form the inclusion of $z$ modulo
$\varpi^2$ and its rows $u$ in $N$ satisfy $su\equiv s\iota_e$ modulo $\varpi^2$ and $u\equiv\iota_e$ modulo $\varpi$; the lattice is \emph{glued for $s$} if $\gamma^{-1}x\equiv(1;0)$ modulo $\varpi$ holds exactly when $xY$ is block normalised for some $Y\in\bar J^1$, then
unique, and after a column operation the columns of $\varpi\gamma^{-1}$ in $V''$ vanish modulo $\varpi$ outside the copy of $a$.

\begin{lemma}[Glued lattices]\label{lem:glued-lattices}
\begin{enumerate}[label=\textup{(\alph*)}]
\item If $Y$ modulo $\varpi^2$ has $\varpi Y=0$, then $1-Y$ is the inverse of an element of $\bar J^1$. Lifts of elements of $\bar J^1$ reduce to $1$ at every chosen place; if an integral point $j$ is $1$ modulo $\varpi^2$, then $\gamma^{-1}j\gamma$ is integral and
reduces to $1$ at every chosen place.
\item If integral $C$, $G$ satisfy $CG=\varpi\cdot1$, some divided lattice has $\varpi\gamma^{-1}\equiv G$ modulo $\varpi^2$.
\item For every $s$ there is a lattice glued for $s$; for a glued lattice the indicator of $(\varpi\gamma^{-1})x\equiv\varpi(1;0)$ is the sum over $Y\in\bar J^1$ of the indicators that $xY$ is block normalised.
\end{enumerate}
\end{lemma}

\begin{proof}
(a) $\varpi Y\equiv0$ says that every entry of $Y$ is divisible by $\varpi$, so $Y^2=0$ modulo $\varpi^2$ and $1-Y$ is a unit fixing $\varpi\cdot1$. A lift of $Y\in\bar J^1$ is $1$ modulo each chosen place by the equation $\varpi Y=\varpi$. For the conjugate, write
$\gamma^{-1}j\gamma-1=\gamma^{-1}(j-1)\gamma$; at a chosen place $\gamma$ and $\varpi\gamma^{-1}$ are integral and $j-1$ is divisible by $\varpi^2$; at the conjugate places use that $\varpi\gamma^{-1}$ and $\bar\varpi\gamma$ are integral everywhere and that $j-1$ is divisible
by $\varpi^2\bar\varpi$, because $j^{-1}=j^*$ is also $1$ modulo $\varpi$; at each place one of $\varpi$, $\bar\varpi$ is $1$. (b) No chosen place is conjugate to a chosen place, so $C$ and $C^*$ commute where needed; put $\varpi\gamma^{-1}$ equal to $G$ at the chosen places,
$\varpi C^*$ at their conjugates and $\varpi$ elsewhere; then $\gamma$ is unitary, by the relation $e\bar e=0$ for the idempotent $e$ of the chosen places. (c) Lift $s$ to an integral $t$, take a Hermite form of $t$ stacked over $\varpi\cdot1$, and reduce:
there are integral $B$, $D$ with $BD=\varpi\cdot1$ such that modulo $\varpi^2$ each of $D$ and the stacked matrix $w$ is a left multiple of the other. The \emph{glue matrix} has the blocks $D$, $-D\iota_e\iota_z^{\mathsf T}$, $0$ and the diagonal matrix that is $\varpi$ on
the range of $z$ and $1$ elsewhere; reading the equation by blocks, $x$ is block normalised exactly when its rows $z$ in $m$ form $1$ and the glue matrix takes $x$ to $\varpi\iota_{z}$. The matrix $C$ with blocks $B$, $\iota_e\iota_z^{\mathsf T}$, $0$ and the diagonal
matrix with $\varpi$ off the range of $z$ satisfies $C\cdot\mathrm{glue}=\varpi\cdot1$ (only the top right block needs an argument, the rows $z$ of the diagonal block being $\varpi$ times those of $1$), so (b) gives a divided lattice with
$\varpi\gamma^{-1}\equiv\mathrm{glue}$, after reordering the basis so that $z$ goes to the indices of $W$. Right multiplication by $Y\in\bar J^1$ fixes $\varpi\cdot1$ and does not change the congruence, and the rows $z$ of $x$ form $Y^{-1}$ exactly when those of $xY$ form
$1$; so $xY$ is block normalised exactly when the rows $z$ of $x$ form $Y^{-1}$ and $(\varpi\gamma^{-1})x\equiv\varpi(1;0)$, which gives existence and uniqueness of $Y$. An invertible $K$ clears the columns of $s$ outside the copy of $a$, and $\varpi$ kills the rows
$\varpi K$, which gives the last condition. The indicator identity is the existence and uniqueness of $Y$.
\end{proof}

\begin{lemma}[Descent of product forms]\label{lem:product-descent}
Let the lattice $\gamma$ be glued for $s$ and let the characters $b$, $b'$ of the two summands have product the character of $V$.
\begin{enumerate}[label=\textup{(\alph*)}]
\item The descent to $W$ of the form supported on the class of $h\gamma$ with value $v$ is, at $g$, the conjugate of $c_\gamma$ times the theta kernel at $(g,h)$ with polynomial the Fock polynomial of $\bar v$ and finite datum the function that is $0$ unless
$(\varpi\gamma^{-1})S\equiv\varpi\iota$ modulo $\varpi^2$ and is then the product over the places of $Q$ of the residue characters at the determinant of the rows of $W$ of $S$.
\item The Fock polynomial of the Cartan product of $y$ (weight $\mu^+$ on $V''$) and $\eta$ (weight $\mu^-$ on $V^\circ$) is the product of the polynomial of $y$ in the positive variables on the rows of $V''$ and that of $\eta$ in the negative variables on the rows of
$V^\circ$.
\item The descent of the product form on the class of $\mathrm{diag}(h,h')\gamma$ is, at $g$, the conjugate of $c_\gamma$ times $\sum_{Y\in\bar J^1}$ of the product of the kernel of $(W,V'')$ at $(g\,\mathrm{lift}(Y),h)$, with the column datum of the rows $w$ and the polynomial of $\bar y$,
and the kernel of $(W,V^\circ)$ at $(g\,\mathrm{lift}(Y),h')$, with the indicator of the inclusion of $z$ and the polynomial of $\bar\eta$; and the same summed over the classes $h$ of the row-fixing level.
\item For a prime $p$, an embedding into $\C_p$ and $M\ne0$ divisible by $\varpi$, the gauge of the Cartan product at $\mathrm{diag}(h,h')\gamma$ is at most $\norm M^{-l_1}$ times the gauge of $y$ at $h$ and that of $\bar\eta$ at $h'$, where the
\emph{denominator exponent} $l_1$ is $\sum_\tau(\#a\cdot\#\{\text{columns of }\mu^+_\tau\text{ below }\#N\}+\#m\cdot\#\{\text{columns of }\mu^-_\tau\text{ below }\#m\})$.
\end{enumerate}
\end{lemma}

\begin{proof}
(a) Split off $\gamma$ from the point $(g,h\gamma)$. Supported on the chosen places, it acts on the finite datum by the scalar $c_\gamma$ and the substitution $y\mapsto\gamma^{-1}y$. The extension of a function on matrices modulo $\varpi^2Q$ to adelic vectors (by $0$
off the integral ones at the places of $\varpi^2Q$) at $\gamma^{-1}x$ equals the extension of the function with support $(\varpi\gamma^{-1})T\equiv\varpi\iota$ at $x$: each side is $0$ unless its vector is integral on the level, and since the square of the class of $\varpi$
vanishes modulo $\varpi^2$, the entries of $\varpi^2\gamma^{-1}x$ vanish, which forces integrality of $\gamma^{-1}x$; by unitarity the conjugate of $\varpi\gamma^{-1}$ is $\bar\varpi\gamma^{\mathsf T}$, integral, as $\bar\varpi$ is $1$ at a place of the level. The point $\gamma$ is
$1$ at the places of $Q$, which do not change $\det\gamma$. (b) The product is read on the image of the mixed Cartan map, the tensor of the included vectors; the Schur matrix of an inclusion matrix acts on polynomials by renaming the rows of
the variables, with no determinant factor, when the diagram has no column of a height between the two ranks. (c) Apply (a) on the left. On the right, join the two kernels by the rule for orthogonal sums and move the lift of $Y$ into the finite
datum; both sides are values of the bilinear form of the kernel, and the finite data and the polynomials agree by (b), \cref{lem:glued-lattices}(c), and the computation of the residue characters: if the rows $S$ of $xY$ in $N$ satisfy
$wS=w\iota$ modulo $\varpi^2$, then at the chosen places the rows $e$ of $S$ have determinant $1$, and at the places of $Q$ the lift of $Y$ reduces to $1$, so the product of the residue characters at the determinant of the rows of $W$ of $x$ is the
character modulo $\varpi^2Q$ of the determinant of the same rows of $x\,\mathrm{lift}(Y)$ divided by the residue characters at $\det\mathrm{lift}(Y)$. Summing over the classes is linear. (d) The Cartan product is an integral matrix applied to the tensor of two Schur
images, one for each block, so its gauge is at most the product of theirs. For each block, lift the relevant change of columns to an integral $k$ and write $h=h'k^{-1}$; the columns of $(g\gamma)^{-1}$ factor as $A\Delta h^{-1}$ with $A$ integral and $\Delta$ diagonal, $M^{-1}$ on a
set $S$ of columns and $1$ elsewhere: $A$ is integral because $\varpi A$ vanishes modulo the chosen places, the columns in $S$ being multiples of $M$ and $\varpi$ times the others vanishing modulo $\varpi^2$. The Schur matrix of $A$ does not increase norms;
that of $\Delta$ is diagonal with entries $M^{-1}$ to the number of entries of a filling in $S$, at most $\#S$ times the number of columns below the height. On $V''$ only the columns in the copy of $a$ cost a power of $M$; on $V^\circ$, read through the
conjugate embedding, every column may.
\end{proof}

\lean{New.Rows.hookNumberMinus, New.Data.FormTuples, New.Rows.IsIntegralFormTuple, New.Rows.formTupleNormSq, New.Rows.card_firstFactorIndex_le_two_mul_n, New.Rows.card_rowIndex_sum_fin_le_two_mul_n, New.Rows.nonempty_carrierIndex_embedding_firstFactorIndex, New.Rows.hasColumnsBelow_of_eq_zero, New.Rows.card_padIndex_sum_flatIndex_sum_secondFactorIndex, New.Rows.hasColumnsBelow_colsPlus, New.Rows.hasColumnsBelow_colsMinus, New.Rows.splittingV, New.Rows.gluingLevel, New.Rows.classAut_gluingLevel_eq_bot, New.Rows.glueLattice, New.Rows.isGlued_splittingV_of_isGlued, New.Rows.isGlued_glueLattice, New.Rows.weilScalar_glueLattice, New.Rows.firstFactorScale, New.Rows.firstFactorScale_ne_zero, New.Rows.firstFactor, New.Rows.firstFactorPairing, New.Rows.firstFactorPairing_eq_mul_thirdExchangeLeftSide, New.Rows.mul_firstFactorPairing_mem_rapidFunctions, New.Rows.exists_mul_firstFactorPairing_eq, New.Rows.twistFactor, New.Rows.exists_norm_twistFactor_sq_eq, New.Rows.twistFactor_mul_firstFactorPairing, New.Rows.isIntegralAndSmall_firstFactor, New.Rows.weylGauge_firstFactor_le_one, New.Rows.weylNormSq_firstFactor_le, New.Rows.inner_adjoint_glue, New.Rows.firstFactorChar_mul_upsilonU1_zpow, New.Rows.goodForm_mul_liftMod_congruenceGroup, New.Rows.conj_normalizedDescent_glue, New.Rows.inner_adjoint_glue_goodLift_eq_integral, New.Rows.mul_mul_thetaKernelPairing_mem_rapidFunctions, New.Rows.inner_adjoint_glue_goodLift_eq_sum_integral, New.Rows.glue_mem_integralForms, New.Rows.colsMinus_le, New.Rows.card_secondFactorIndex_le, New.Rows.sum_mul_colsMinus_le, New.Rows.hookNumberMinus_le, New.Rows.card_columnFilling_colsMinus_le, New.Rows.card_stdClasses_le_exp_of_le, New.Rows.principalLevel_le_gluingLevel, New.Rows.M_pow_denominatorExponent_sq_mul_le_exp, New.Rows.weylNormSq_evalForm_productForm_le, New.Rows.piM_sq_mem_placeAt_descentPair_one, New.Rows.normSqAt_glue_le, New.Rows.hyperplaneSum_apply_eq_zero, New.Rows.normalizedDescent_eq_zero_of_mem_orthogonal, New.Rows.tamePart_twistPlace_ne_one, New.Rows.stdClassSum_gramForm_eq_zero_of_mem_falseLiftsAt, New.Rows.sum_slice_eq_zero_of_mem_falseLiftsAt, New.Rows.thetaKernelPairing_adjoint_glue_eq_zero, New.Rows.glueFactor, New.Rows.glueFactor_ne_zero, New.Rows.twistFactor_mul_inner_adjoint_glue, New.Rows.periodLift, New.Rows.inner_periodLift, New.Rows.twistFactor_smul_adjoint_glue_comp, New.Rows.isOrthogonalToFalseLifts_glue_periodLift, New.Rows.periodLift_injective, New.Rows.twistFactor_ne_zero}

\begin{theorem}[Gluing]\label{thm:gluing}
\begin{enumerate}[label=\textup{(\alph*)}]
\item (Seesaw.) $c'(\vartheta)\,\mathrm{glue}_\vartheta^*\Theta_{j,\vartheta}=\Theta^\circ$ on good forms, and $\langle\Theta^\circ y,\zeta\rangle=\sum_ic^\circ_i\int y\,\mathsf F_i\,\langle\zeta_i,\theta^\circ(g,\cdot)\rangle\,dg$, with $\theta^\circ$ the theta kernel of $(W_j,V^\circ)$.
\item $\Theta^\circ$ is injective.
\item $\mathrm{glue}_\vartheta$ keeps integrality at every prime not dividing $\Nm Q$, and $\norm{\mathrm{glue}_1\zeta}_\sigma^2\le\exp(2fmb\log^2b)\norm\zeta_\sigma^2$.
\item The adjoint of the gluing takes a false lift to a tuple whose pairing with $\theta^\circ(g,\cdot)$ vanishes for every $g$; and $\mathrm{glue}_\vartheta(\Theta^\circ y)$ is orthogonal to the false lifts.
\end{enumerate}
\end{theorem}

Here $\mathrm{glue}_\vartheta\zeta$, for $\zeta$ in the copy $i$, is $M^{l_1}$ times the sum over the classes $h$ of level $J_\Lambda$ and $h^\circ$ of level $J(\varpi^2)$ of the product forms on the classes of $(h,h^\circ)\mathsf g$, and $c^\circ_i=M^{l_1}c_{\mathsf g}\#\bar J^1c''$.
The level $J(\varpi^2)$ on $V^\circ$ and the level $J_\Lambda$ on $V''$ (the integral points that are $1$ modulo $\varpi$, fix $\Lambda$ modulo $\varpi^2$, and modulo $\varpi^2Q$ have the columns of $1$ off $W_j$) are free: all stabilisers of classes are trivial
(\cref{prop:free-levels}); so a tuple is the table of its values at representatives and inner products are plain sums. The lattice $\mathsf g$ is chosen once, at the trivial twist; the conditions on a glued lattice do not involve the twist, and its
scalar does not depend on it.

\begin{proof}
(a) \emph{The period as a sum of slices.} By the definition of the gluing, $\langle\zeta,\mathrm{glue}_\vartheta^*\varphi\rangle$ is $M^{l_1}$ times the sum over the choices $i$, the classes $c$ of $\U(V^\circ)$ with representatives $h^\circ$, and the classes $h$ of level $J_\Lambda$ of the
pairing of $\varphi''_{\vartheta,i}(h)\cdot\zeta(i,c)$ with $\varphi((h,h^\circ)\mathsf g)$. \emph{On a lift.} Let $\varphi=\Theta_{j,\vartheta}y$. The lift is $\int y(g)\theta(g,\cdot)\,dg$, and $\langle\theta(g,\cdot),\mathrm{glue}\,\zeta\rangle$ is the descent of $\mathrm{glue}\,\zeta$ (\cref{lem:vector-kernel}(c)),
so the pairing is the integral of $y$ against the conjugate of the descent. By \cref{lem:product-descent}(c) the conjugate of the descent at $g$ is the sum over $i$ and $u\in\bar J^1$ of $M^{l_1}c_{\mathsf g}\,\Xi_{\vartheta,i}(gu)\langle\zeta_i,\theta^\circ(gu,\cdot)\rangle$: the
characters satisfy $\nu''\cdot\upsilon^{m/2-j}=$ the character of $V_j$, because $m$ is even and $k-1=r-m$. A good form is invariant under the lifts of the elements of $\bar J^1$, which lie in the level $J(\varpi)$; the measure is right
invariant, so all $u$ give the same integral, and the products are rapidly decreasing (finite combinations of kernels times a form). Hence
\[
  \langle\zeta,\mathrm{glue}_\vartheta^*\Theta_{j,\vartheta}y\rangle=\sum_iM^{l_1}c_{\mathsf g}\#\bar J^1\int y\,\Xi_{\vartheta,i}\,\overline{\langle\zeta_i,\theta^\circ(g,\cdot)\rangle}\,dg.
\]
Multiply by $c'(\vartheta)$ and use \cref{prop:first-factors}(c): $c'(\vartheta)\Xi_{\vartheta,i}=c''\mathsf F_i$, free of the twist. The right side is then $\langle\zeta,\Theta^\circ y\rangle$ with $\Theta^\circ$ defined at the trivial twist, which proves both claims.

(b) If $\Theta^\circ y=0$, test against the tuple that is $\zeta$ in the copy $i$ and $0$ elsewhere: the copies other than $i$ contribute nothing, and what is left is $c^\circ_i\ne0$ times $\int y\,\mathsf F_i\,\langle\zeta,\theta^\circ(g,\cdot)\rangle\,dg$, which is therefore $0$ for all
$i$ and $\zeta$. \Cref{thm:products-span} gives $y=0$. As a nonzero good form exists, $c'(\vartheta)\ne0$.

(c) \emph{Integrality.} All stabilisers are trivial, so $\mathrm{glue}\,\zeta$ is $M^{l_1}$ times a plain sum of forms supported on one class, and the gauge of such a form is that of its vector $\varphi''(h)\cdot\zeta(h^\circ)$; by \cref{lem:product-descent}(d) this is at most
$\norm M^{-l_1}$ times the gauges of the two factors, which are integral by \cref{prop:first-factors}(b) and the hypothesis. \emph{Norm.} At the trivial twist the level $J_\Lambda$ contains $J(\varpi^2)$, so the numbers of classes of $\U(V_j)$ and $\U(V'')$ are at
most $\exp b$ (levels containing $J(\varpi^2)$ in rank at most $2n_m$). A form supported on one class costs the order of a stabiliser, at most $n_{\rm Aut}$, and the norm of the product vector is the product of the norms, the first at most $\exp(fmb\log^2b)$.
The norm of a sum is at most the number of terms times the bound, and $(M^{l_1})^2\exp(b)(n\exp(b)n_{\rm Aut})^2\le\exp(fmb\log^2b)$, where $n$ counts the pairs of a choice and a class, by a comparison of the sizes of the data, in which $l_1$ is $3m$ times the
number of columns of $\mu^+$ plus $\operatorname{rk}V^\circ$ times the number of columns of $\mu^-$, summed over the places, and the numbers of matrices, classes and selections of rows are counted.

(d) \emph{Periods of false lifts.} Let $\Xi$ be a false lift: in the image of $e_0$, killed by the Hecke polynomial operator at a split place off the level, and orthogonal to the lifts of good forms. Its descent vanishes: for every $y$,
$\langle\Theta y,\Xi\rangle=\langle\Theta y,e_0\Xi\rangle=\langle e_0\Theta y,\Xi\rangle=0$, as $e_0\Theta y$ is the lift of a good form. By \cref{thm:false-lift-periods}, the class sum of the pairing of $C(k(h))$ with $\Xi((h_1,h)x)$ vanishes for every conjugate-linear equivariant $C$
and every function $k$ whose coordinates are multiples of conjugate theta kernels of $(W_j,V^\circ)$; take $C(\eta)=y\cdot\eta$ for $y$ of weight $\mu^+$, the point $(h,h^\circ)\mathsf g$, and the level $J(\varpi^2)$, which fixes the pairing because its conjugate by $\mathsf g$ fixes
$\Xi$ (\cref{lem:glued-lattices}(a)); the level is free, so the class sum is a plain sum, and the inner product with a form supported on one class is the pairing with the value there. Hence the slices of $\Xi$ vanish, and the conjugate of the
value of $\mathrm{glue}^*\Xi$, paired with the tuple of values of $\theta^\circ(g,\cdot)$, is a sum of slices: the kernel is linear in its polynomial, so expanding both in the basis of weight $\mu^-$ gives $0$. \emph{Orthogonality.} For a false lift $\Xi$,
$\langle\Xi,\mathrm{glue}\,\Theta^\circ y\rangle=\langle\mathrm{glue}^*\Xi,\Theta^\circ y\rangle$, which by (a) is an integral against the products of the pairings of $\mathrm{glue}^*\Xi$ with the kernels $\theta^\circ(g,\cdot)$, and these vanish.
\end{proof}

\subsection{Products span}\label{ss:products-span}
For a real point $h$ let $\abs h$ be its largest operator norm and $\lambda_1(g)$ the first minimum of the lattice at $g$, cut off at $1$. A \emph{peak} of a bounded function is a point where its modulus exceeds half the
supremum.
\lean{New.Data.negIndex_posRowIndex_injective_and_padSigns, New.Data.exists_datum_norm_modKernel_sub_polyMatrixCoeff_le, New.Data.paddedMinorProd_mem_supported, New.Data.maxOpNorm_sum_one_le, New.Data.polyMatrixCoeff_paddedMinorProd_ne_zero_and_bounded, New.Data.exists_eval_smul_paddedMinorProd, New.Data.flatKernelWeight, New.Data.holomorphicWeight, New.Data.antiholomorphicWeight, New.Data.polyMatrixCoeff_paddedMinorProd, New.Data.flatKernelWeight_antiholomorphicWeight_bounds, New.Data.carrierChar_hasWeight_sub_card, New.Data.exists_isPeak_carrierForm, New.Data.paddedSpace, New.Data.firstMinimum_le_paddedSpace, New.Data.exists_sum_choiceFn_eq_integral, New.Data.inv_pow_twenty_mul_pow_nStar, New.Data.inv_pow_le_firstMinimum_mul, New.Data.card_carrierIndex_le_and_card_sum_padIndex_le, New.Data.integral_mul_polyMatrixCoeff_paddedMinorProd, New.Data.exists_choiceFn_approx, New.Data.thetaKernelPairing, New.Data.isAdelicHol_star_goodForms, New.Data.integral_mul_sum_choiceFn_eq_zero, New.Data.eq_zero_of_integral_choiceFn_thetaKernelPairing}
\lean{New.Fock.norm_sq_le_norm_sq_add_norm_mulVec_add_sq, New.Fock.norm_fockInner_realHeisenbergOp_le, New.Rows.firstMinimum_mul_le_hom, New.Rows.dataInner_traceChar, New.Rows.norm_pairHomInf_le_mul, New.Rows.le_norm_realCoord, New.Rows.le_norm_realCoord_sub, New.Rows.tsum_inv_norm_realCoord_pow_le, New.Rows.sum_dataInner_traceChar_div_card, New.Rows.modKernel_expect_traceChar, New.Rows.exists_norm_modKernel_sub_fockInner_le, New.Data.eight_mul_pow_mul_two_pow_le, New.Data.two_pow_mul_pow_f_le_b_pow_nStar, New.Data.two_mul_card_pow_three_le_nStar, New.Data.eight_mul_card_sq_lt, New.Data.hasLargeGap_and_dsCoeff_eq_holMinors, New.Data.integrable_tail_and_integral_le, New.Data.norm_holMinorsPi_mul_pow_le, New.Data.integrable_norm_holMinorsPi_mul_pow, New.Data.integral_norm_holMinorsPi_sq_eq_inv_prod, New.Data.integral_norm_holMinorsPi_sq_pos, New.Data.integral_norm_holMinorsPi_mul_pow_le, New.Data.integrable_mul_holMinorsPi, New.Data.integrable_holMinorsPi, New.Data.norm_integral_mul_holMinorsPi_mul_le, New.Data.integral_mul_norm_apply_one_mul_le, New.Data.norm_apply_one_le_of_forall_isPeak}

The archimedean tools concern the groups $\U(\mathrm{diag}\,s)$ of rank $1\le n\le t$, the \emph{basic holomorphic function} $m_u$ of a weight $u$ (a product of quotients of principal minors, one factor for each place), and orders $o$ of the indices that put the
signs $+1$ first. A weight is \emph{$K$-dominant} for $o$ if it decreases weakly along $o$ within each sign, and has \emph{gap} $g$ if its entries at the signs $-1$ exceed those at the signs $+1$ by at least $g$. A function on the real group is
\emph{holomorphic of weight $u$} if it extends analytically to products with $o$-triangular matrices, on which it transforms by the character $u$.

\begin{lemma}[Fock coefficients and sparse sums]\label{lem:sparse-sums}
\begin{enumerate}[label=\textup{(\alph*)}]
\item Let $P$ be a homogeneous polynomial of degree $d$ with Fock vector $v$, let $\sigma$ be unitary, $y\ne0$, and let $A\ge d+\#\iota$ and $n$ be natural numbers. Then
\[
  \abs{\langle v,\varrho(y)\omega(\sigma)v\rangle}\le(4(A+n))^n\abs y^{-2n}\norm\sigma^{\#\iota+2n}\langle v,v\rangle.
\]
\item If $N\ne0$ lies in the different of $E$, then $\lambda(S)\lambda(T)/(C\#n'\#n)\le\lambda(\Hom(S,T))$ with $C=6N(f\#n)^2$.
\item Let $c\ne0$ and $m>0$ satisfy $m\abs{c}_w\le\sqrt{2\pi\im\jmath_w}\,\lambda$ at every infinite place. If $y\ne0$ and $cy$ is a lattice vector, the Fock coordinates of $y$ have norm at least $m$; distinct such vectors are at
distance at least $m$; and for $q$ above the number of coordinates, $\sum\abs y^{-2q}\le(5/m)^{2q}$.
\item (Sparse sums.) Let $V$, $V'$ have ranks at most $t$, $Q$ homogeneous of degree $d\le b^{3/2}$, and $\lambda_1(g)\lambda_1(g')\ge b^{-13}$. Then for some datum $\Psi$ modulo $\varpi^2$ the theta kernel with $\Psi$, $Q$ at $(gh,g'h')$ differs from the Fock
coefficient $c_Q(h,h')$ by at most $b^{-20n_*}(\abs h\abs{h'})^{6n_*}\abs{c_Q(1,1)}$, for all real $h$, $h'$.
\end{enumerate}
\end{lemma}

\begin{proof}
(a) The coefficient is at most $\norm\sigma^{\#\iota}\int\norm{v(x)}\norm{v(\sigma x+y)}\,dx$. Since $\abs y^2\le2\norm\sigma^2(\abs x^2+\abs{\sigma x+y}^2)$, the points $x$ and $\sigma x+y$ are not both small. Dilating by $c=\sqrt{n/A+1}$ takes a Gaussian factor out
of $v$, which is small at $x$ or at $\sigma x+y$; by homogeneity of $P$, $1-c^{-2}=n/(n+A)$, and undoing the dilation costs $c^{2(d+\#\iota)}\le e^n$. The rest is bounded by $2ab\le a^2+b^2$, both terms having the same integral, as $x\mapsto\sigma x+y$
preserves the measure; by the reproducing property of the vacuum the integrand has the norm bounded. (b) Read a nonzero lattice vector $z$ of $\Hom(S,T)$ as a matrix and apply it to a short basis of the lattice of $S$ (\cref{lem:deep-poincare}
uses the same comparison): some image is a nonzero lattice vector of $T$, of length at least $\lambda(T)$ and at most $\#n'\#n\norm z$ times the length of the basis vector. (c) $cy$ is a nonzero lattice vector, of length at least the first minimum;
compare coordinates; for distinct vectors apply this to the difference; the sum is the packing estimate of \cref{lem:deep-poincare}(c) with exponent $2q$ above the real dimension by at least $2$. (d) Take for $\Psi$ the mean of the trace characters
$\psi_t=\psi(\tr(\cdot\,t))$ weighted by the conjugates of their theta kernels at $(g,h)$; the unit scalars of the characters cancel against their conjugates, and the real points act on the Fock vector by the real Weil representation. The
kernel of $\Psi$ at $(gu,hu')$ is then the mean over $t$ of products of thetas, which by the computation of the pairing of trace data (the finite part of the datum of $\psi_t$ is the vacuum translated by a finite vector; moving
Heisenberg operators past it costs a phase, and the sum over $t$ is the indicator that the finite coordinates vanish modulo the modulus) is a sum over rational vectors $y$ of the lattice scaled by $c=q_1/\varpi^2$ of terms
$\langle v,\varrho(\xi_y)\Omega(u,u')v\rangle$ times a phase. The term $y=0$ is the Fock coefficient. The scaled lattice points are $b^6$-separated: the first minimum of $\Hom(S,T)$ is at least $\lambda_1(g)\lambda_1(g')$ over a constant by (b), and $\abs c$
is small. Each other term decays by (a), with $n=2n_*$; summing by (c) and collecting the constants ($100=4\cdot5^2$) gives $(100(d+2q)/m^2)^q(\abs u\abs{u'})^{3q}\abs{\langle v,v\rangle}$, which is at most the stated bound for the sizes of the data.
\end{proof}

\begin{lemma}[Holomorphic coefficients]\label{lem:holomorphic-integrals}
Let $1\le n\le t$, $o$ put the signs $+1$ first, and $u$, $u'$ be dominant with gap $300n_*$, $u$ with entries at most $5b$.
\begin{enumerate}[label=\textup{(\alph*)}]
\item $l=2v+2\rho$ (with $2\rho$ the signed count for $-o$) is a doubled parameter of gap $500n^3+8n_*$, and its discrete series coefficient is the quotient $m_{v+c}/m_c$ of products of principal minors, $c_i=2\sum\abs{v_j}$ at a sign $+1$ and $0$
otherwise.
\item $\int\abs{m_u}^2$ is the inverse of the product of the formal degrees, hence positive; $\abs{m_{u'}(h)}\abs h^{6n_*}\le2^{6n_*}\prod_\tau(2/\norm{h_\tau})^{10n^3}$, which is integrable with integral at most $b^{n_*}\int\abs{m_u}^2$; and $m_u$ is integrable.
\item If $A$ is continuous with $\abs{A(h)}\le C\abs h^{6n_*}$, then $\abs{\int Am_{u'}}\le Cb^{n_*}\int\abs{m_u}^2$. If $W$ is bounded by $M$, and $K$, $K'$ are within $\delta\abs h^{6n_*}$ of $K_0$ with $\int WK'm_{u'}=0$, then
$\abs{\int Km_{u'}W}\le2\delta Mb^{n_*}\int\abs{m_u}^2$.
\item (Isolation.) Let $F$ on the adelic group be continuous, holomorphic of weight $u$ in the real variable, with a peak, and $K_0$ bounded. If at every peak $x$, $\int\overline{F(xh)}K(h)m_{u'}(h)\,dh=0$ for some $K$ holomorphic of weight $u-u'$ with
$\abs{K-K_0}\le\delta\abs h^{6n_*}$, then $\abs{K_0(1)}\le(4b^{n_*}+1)\delta$.
\end{enumerate}
\end{lemma}

\begin{proof}
(a) Dominance within the blocks and the gap make the entries of $l$ distinct with the right parity and the stated gap; the coefficient of the holomorphic discrete series is then the quotient of minors (\cref{ss:A-coefficients}). (b) At each
place the factors of $m_u$ and $m_{u'}$ are discrete series coefficients, so $\int\abs{m_u}^2$ is the inverse product of formal degrees by Schur orthogonality. For the bound, use $6n_*$ more powers of decay at a place of largest norm: the tail
$\deg(l)(2/\norm g)^{10n^3}$ is integrable on $\U(\mathrm{diag}\,s)$ with integral at most $2^{13t^3}b^{14t^2}$, as $10n^3>8n^2$ is above the threshold for integrability and the constant of the tail bound is at most this for $L\le10b+t$; and
$2^{6n_*}(2^{13t^3}b^{14t^2})^f\le b^{n_*}$, with the formal degrees those of $u$. (c) The first is (b); for the second, orthogonality lets one replace $K$ by the small function $K-K'$. (d) Rotate the peak $z$ by a circle $\mathrm{rot}(\theta)$: all $z\,\mathrm{rot}(\theta)$
are peaks, so the hypothesis gives functions $K_\theta$, all within $\delta\abs h^{6n_*}$ of $K_0$. Replacing $K_\theta$ by the function $K_{(0)}$ of the angle $0$ costs, by the second part of (c), at
most $2\delta\sup\abs F\,b^{n_*}\int\abs{m_u}^2$ in each integral. Let $W_\theta$ be the conjugate of the quotient of the rotated translate of $F$ by the character; as $h\mapsto F(zh)$ is
holomorphic of weight $u$, the circle mean of $W_\theta$ is $\overline{F(z)m_u}$. Exchanging the integrals over $\theta$ and $h$ and using the reproducing formula for the holomorphic function $K_{(0)}m_{u'}$ of weight $u$, whose value at $1$ is $K_{(0)}(1)$, gives
$\int\abs{m_u}^2\abs{F(z)}\abs{K_{(0)}(1)}\le2\delta\sup\abs F\,b^{n_*}\int\abs{m_u}^2$. As $\sup\abs F<2\abs{F(z)}$, $\abs{K_{(0)}(1)}\le4b^{n_*}\delta$, and $\abs{K_{(0)}(1)-K_0(1)}\le\delta$.
\end{proof}

\begin{lemma}[Weights on $W_j$ and the carrier]\label{lem:weights-Wj}
\begin{enumerate}[label=\textup{(\alph*)}]
\item The weight $u$ of the conjugates of good forms (the highest weight of the lowest $K$-type of $-2\varsigma$, for the target parameter $\varsigma$ of \cref{ss:grids}) is the sum of two staircase weights: one on $V''$ with the columns of $\mu^+$ in sign $+1$ and shift $\kappa_j-j-\lfloor(t+k-1)/2\rfloor$, and the
\emph{kernel weight} $u^\circ$ on $V^\circ$ with the columns of $\mu^-$ in sign $-1$ and shift $j-\lfloor m/2\rfloor$. Both $u$ and $u^\circ$ are $K$-dominant for the order $o_W$ that lists the positive vectors first, in reversed order, and have gap $300n_*$; the
entries of $u$ are at most $5b$.
\item The functions $\mathsf F_i$ are holomorphic of weight $u-u^\circ$; good forms have full level at every multiple of $\varpi$ and are rapidly decreasing; the conjugate of a good form is holomorphic of weight $u$; and $\mathsf F_i$ is a multiple of the
theta series of $(W_j,V'')$ at the trivial twist, so $y\mathsf F_i$ is rapidly decreasing of level $\varpi^2$.
\item The matrix coefficient $c_P$ of the Fock vector of the padded polynomial on $(S,\tilde W)$ satisfies $c_P(h,1)=c_P(1,1)m_2(h)\overline{m_3(h)}$, with $m_2$, $m_3$ the basic functions of $(\Delta';\#\mathrm{Det}+r+2m)$ and $(0;\#\mathrm{Det}+r+11m+1)$; $c_P(1,1)\ne0$ and $c_P$ is bounded;
$P$ is a polynomial in the variables of sign $1$ of degree at most $b^{3/2}$. The flat weight $u_1=(-\Delta';6m+1)$ and $u_3$ are dominant with gap $300n_*$ for every order, $u_3$ with entries at most $5b$.
\item The carrier has a peak $g_1$ with $\lambda_1(g_1)\ge(5b\sqrt b)^{-1/2}$.
\end{enumerate}
\end{lemma}

\begin{proof}
(a) By the formula for the lowest $K$-type as a staircase weight, $u$ is one staircase weight for $V_j=I_{r+t}$, whose sum runs up to $r+t=6r$; the sums on the right stop at the ranks of $V''$ and $V^\circ$, and the missing terms vanish. On
positive vectors $u$ is a constant minus a staircase sum and on negative vectors a constant plus $u^\circ$; each staircase sum decreases as its starting point grows, which gives dominance. At a negative vector the sum of the columns of
$\mu^-$ above its number is at least $300n_*$ (the columns of $\mu^-$ are at most $\Delta$ in each height and the heights fill the range), and the difference $u-u^\circ$ is larger at negative than at positive vectors; the bound on the entries is the
bound on the parameter plus $r$. (b) $\mathsf F_i$ is a finite combination of kernels of $(W_j,V'')$ with polynomials in the positive letters (\cref{prop:three-exchanges}(c) at the trivial twist), the definite group acts on these polynomials by a
linear substitution, and the Fock vector of the highest vector is holomorphic of weight $u-u^\circ$, so all are. A good form is a descent (the lift is injective, so its adjoint composed with it is onto in finite dimension), the polynomial of the
highest vector is a product of joint harmonics, hence holomorphic, and by equivariance so is the polynomial of every vector; the conjugate of the descent is a finite combination of theta kernels with such polynomials. The tame part of the
characters is trivial off the modulus. (c) Each minor of $P$ uses all the columns of one sign, so a substitution multiplies it by a principal minor; entry by entry, the weights are as stated, the character of the kernel having weight
$m-\#\mathrm{Det}$ (\cref{lem:exchange-numerics}(d)). The Fock representation is unitary, so $c_P$ is bounded, and $c_P(1,1)=\langle v,v\rangle\ne0$. Each variable of $P$ has row and column of the same sign. (d) The carrier is automorphic of every
level, rapidly decreasing and holomorphic (\cref{prop:three-exchanges}(a)); at the level $\sqrt{-d}$, of norm $\sqrt d\le\sqrt b$, the cusp estimate (\cref{thm:B-cusp}(c)) bounds a form high in a cusp by half its supremum, so a peak has first minimum at
least the stated threshold.
\end{proof}

\begin{theorem}[Products span]\label{thm:products-span}
A good form $y$ on $W_j$ with $\int y(g)\mathsf F_i(g)\theta^\circ(g,\cdot)\allowbreak\,dg=\nobreak0$ for all choices $i$ is zero.
\end{theorem}

\begin{proof}
Suppose $y\ne0$, and write $\tilde W=W_j\oplus I_j^-\oplus I_{3m-j}$.

\emph{Step 1: integrals of the carrier are combinations of the $\mathsf F_i$.} Let $g_1$ be a point of $\U(S)$, $\Psi$ a function modulo $\varpi^2$ and $\theta[\Psi,P]$ the kernel of $(S,\tilde W)$ with the characters $\upsilon^m$ and $\nu^\flat\upsilon^m$. As a
function of $g\in\U(W_j)$, the integral of $\varphi^\circ(g_1h)\theta[\Psi,P](g_1h,(g,1))m_{u_1}(h)$ over the real points $h$ is a combination of the $\mathsf F_i$. Indeed, the reproducing formula on $(S,W^\flat)$ turns the real integral against $m_{u_1}$ into pairings
with kernels of $(S,W^\flat)$: the translates of the power of one minor lie in the span of the products of determinants formed with the short bases of the classes. The product of $\theta[\Psi,P]$ with a kernel of $(S,W^\flat)$ is a kernel of
$(S,\tilde W\oplus W^\flat)$, linear in its datum, and the product of $\Psi$ with the datum is a combination of additive characters, which are the data defining $\mathsf F_i$; the products with $\varphi^\circ$ are rapidly decreasing of level $\varpi^2$.

\emph{Step 2: the hypothesis on the real group.} For every combination $\sum c_i\mathsf F_i$ and every adelic $g$, $\int y(gh)\sum_ic_i\mathsf F_i(gh)\,m_{u^\circ}(h)\,dh=0$. By the reproducing formula on $(W_j,V^\circ)$ for the basis polynomials $Q$ in the negative letters,
whose span holds the translates of the highest vector, the real integral against $m_{u^\circ}$ is a combination of pairings with kernels; the pairing of $y\sum c_i\mathsf F_i$ with a kernel is a combination of the pairings of the $y\mathsf F_i$, and each of these is a
case of the hypothesis for a tuple supported at one class whose value there is the conjugate of a vector with the basis polynomial.

\emph{Step 3: approximation away from the cusps.} Let $g_1$ be the peak of \cref{lem:weights-Wj}(d) and put $F_0(h)=\int\varphi^\circ(g_1h')c_P(h',h)m_{u_1}(h')\,dh'$, a bounded function on the real group of $W_j$. By \cref{lem:weights-Wj}(c), the holomorphic factor of $c_P(\cdot,1)$ times
$h\mapsto\varphi^\circ(g_1h)m_{u_1}(h)$ is holomorphic of the weight of $m_3$, and the reproducing formula gives $F_0(1)=c_P(1,1)\varphi^\circ(g_1)J$ with $J=\int\abs{m_{u_3}}^2>0$ (\cref{lem:holomorphic-integrals}(b)), so $F_0(1)\ne0$. If $\lambda_1(g)\ge(5b^{24})^{-1/2}$, then $\lambda_1(g_1)\lambda_1(g)\ge b^{-13}$ (extending $g$ by $1$ to the padding does not lower the first minimum, both summands being
definite), so by \cref{lem:sparse-sums}(d) some kernel $\theta[\Psi,P]$ is within $b^{-20n_*}(\abs{h'}\abs h)^{6n_*}\abs{c_P(1,1)}$ of $c_P$. By Step 1 the same integral with this kernel is a combination $\mathsf F$ of the $\mathsf F_i$, and by \cref{lem:holomorphic-integrals}(c) the
difference has the same factors $\abs{\varphi^\circ(g_1)}$ and $J$, so
\[
  \abs{\mathsf F(gh)-F_0(h)}\le2b^{-19n_*}\abs{F_0(1)}\abs h^{6n_*},
\]
using $b^{-20n_*}b^{n_*}=b^{-19n_*}$.

\emph{Step 4: isolation.} The conjugate $\bar y$ is holomorphic of weight $u$ (\cref{lem:weights-Wj}(b)) and has level $M\sqrt{-d}$, so by the cusp estimate its peaks have $\lambda_1\ge(5b^{24})^{-1/2}$. At every peak $x$, Step 3 gives a combination
$\mathsf F$, holomorphic of weight $u-u^\circ$, with $h\mapsto\mathsf F(xh)$ within $\delta\abs h^{6n_*}$ of $F_0$, $\delta=2b^{-19n_*}\abs{F_0(1)}$, and Step 2 gives $\int y(xh)\mathsf F(xh)m_{u^\circ}(h)\,dh=0$. By \cref{lem:holomorphic-integrals}(d) for $\bar y$, with $u'=u^\circ$,
$\abs{F_0(1)}\le(4b^{n_*}+1)2b^{-19n_*}\abs{F_0(1)}$. With $B=b^{n_*}>10$, $(4B+1)2<B^2\le B^{19}$, so $\abs{F_0(1)}<\abs{F_0(1)}$, which is absurd.
\end{proof}

\subsection{The comparison operator}\label{ss:comparison-operator}
Let $n_{\rm Aut}$ be a common multiple of the orders of all stabilisers in the class sums and $h_k$ the hook number of the diagram $\mu^-$ of the forms on $V^\circ$.

\begin{theorem}[Comparison operator]\label{thm:comparison-operator}
For $1\le j\le m$ the operators $U_\vartheta=n_{\rm Aut}h_k^2\,\mathrm{glue}_\vartheta\mathrm{glue}_\vartheta^*$ on the forms on $V_j$, one for each twist, satisfy:
\begin{enumerate}[label=\textup{(\alph*)}]
\item $\langle\Theta_\vartheta y',U_\vartheta\Theta_\vartheta y\rangle=\rho_\vartheta\langle y',U^\circ y\rangle$ on good forms, with $\rho_\vartheta=\Gamma_j(\vartheta)/\Gamma_j(1)\ne0$ and $U^\circ$ positive definite and independent of $\vartheta$;
\item $U_\vartheta$ takes lifts to forms orthogonal to the false lifts;
\item $U_\vartheta$ keeps integrality above every prime not dividing $\Nm Q$;
\item at every automorphism $\sigma$ of $\C$, $U_1$ multiplies squared norms by at most $\exp(6fmb\log^2b)$.
\end{enumerate}
\end{theorem}
\lean{New.Data.normDet_sq_mul_det_compression, New.Data.norm_map_det_le_pow, New.LevelPlaces.LevelForm.normSqAt_smul, New.LevelPlaces.LevelForm.normSqAt_nonneg, New.LevelPlaces.LevelForm.normSqAt_pos, New.DescentPair.exists_smul_mem_integralForms_maximal, New.Data.IsComparisonOperator, New.Data.comparisonOp, New.Data.heckeValueProd_div_eq, New.Data.adjoint_glue_goodLift, New.Data.inner_goodLift_comparisonOp, New.Data.isOrthogonalToFalseLifts_comparisonOp, New.Data.TightFrameIndex, New.Data.exists_frame_weightRep_eq_intCast, New.Data.exists_tight_frame_formTuples, New.Data.comparisonOp_eq_sum, New.Data.comparisonOp_mem_integralForms, New.Data.card_prod_mul_card_tightFrameIndex_mul_sq_le, New.Data.normSqAt_comparisonOp_le, New.Data.isComparisonOperator_comparisonOp, New.Data.finrank_range_goodLift, New.Data.range_goodLift_le_range_inf_ker, New.Data.normSqAt_starProjection_range_goodLift_le, New.Data.liftCompression, New.Data.IsComparisonOperator.exists_normDet_sq_mul_det_eq, New.Data.IsComparisonOperator.norm_padic_det_liftCompression_le, New.Data.IsComparisonOperator.norm_map_det_liftCompression_le, New.Data.mFoldTwist_one_pow, New.Data.heightLe_liftDeterminant_ratio}

\begin{lemma}[Tight frames]\label{lem:tight-frame}
\begin{enumerate}[label=\textup{(\alph*)}]
\item For each rational unitary $\gamma$ of $V^\circ$ there is a frame of vectors $\eta_T$ of weight $\mu^-$, one for each filling $T$ with increasing columns, with $\sum_T\langle\eta_T,\xi\rangle\eta_T=h_k^2\xi$ for all $\xi$, $\gamma$ taking the conjugate of each $\eta_T$ to an integer
vector, and each of squared norm at most $h_k^2$ after every automorphism of $\C$.
\item For rational points $\delta_c$, one for each class, there are tuples $f_a$, indexed by a choice, a class and a filling, with $\sum_a\langle f_a,\zeta\rangle f_a=h_k^2\zeta$, $\delta_c$ taking the conjugate of the value of $f_a$ at $(i,c)$ to an integer vector, and squared norm
at most $h_k^2$ at every automorphism.
\item $U_\vartheta\varphi=\sum_an_{\rm Aut}\langle\varphi'_a,\varphi\rangle\varphi'_a$ with $\varphi'_a=\mathrm{glue}_\vartheta f_a$; and the number of pairs of a choice and a class times $(\#\text{frame}\cdot n_{\rm Aut}h_k^2)^2$ is at most $\exp(2fmb\log^2b)$.
\end{enumerate}
\end{lemma}

\begin{proof}
(a) Move an integral tight frame of the Weyl module of $\mu^-$ (vectors $\eta^0_T$ indexed by the fillings with increasing columns, with $\sum_T\langle\eta^0_T,\xi\rangle\eta^0_T=h_k^2\xi$,
integer coordinates and squared norms at most $h_k^2$) by the unitary map $\xi\mapsto\overline{\gamma^{-1}\bar\xi}$. (b)
Put $f_{(b,T)}$ equal to the vector $\eta_{b,T}$ of the frame moved by $\delta_c$ at $b=(i,c)$ and $0$ elsewhere; the metric on the values is the form of the Gram matrix, and the level is free, so inner products are sums over the pairs.
(c) $\mathrm{glue}\,\mathrm{glue}^*\varphi=h_k^{-2}\sum_a\langle\mathrm{glue}f_a,\varphi\rangle\mathrm{glue}f_a$ by the frame identity applied to $\zeta=\mathrm{glue}^*\varphi$. The count is \cref{thm:gluing}(c) with the further factors, all at least $1$.
\end{proof}

\begin{proof}[Proof of \cref{thm:comparison-operator}]
(a) By \cref{thm:gluing}(a), $\mathrm{glue}_\vartheta^*\Theta_\vartheta y=c'(\vartheta)^{-1}\Theta^\circ y$, so
\[
  \langle\Theta_\vartheta y',U_\vartheta\Theta_\vartheta y\rangle=n_{\rm Aut}h_k^2\abs{c'(\vartheta)}^{-2}\langle\Theta^\circ y',\Theta^\circ y\rangle.
\]
By \cref{prop:first-factors}(c), $\abs{c'(\vartheta)}^2=C/\Gamma_j(\vartheta)$ with $C\ne0$, as $c'(\vartheta)\ne0$; hence
\[
  \abs{c'(1)}^2/\abs{c'(\vartheta)}^2=\Gamma_j(\vartheta)/\Gamma_j(1)\ne0.
\] Put $U^\circ=n_{\rm Aut}h_k^2\abs{c'(1)}^{-2}\Theta^{\circ*}\Theta^\circ$; it is independent of $\vartheta$ and positive definite,
$\Theta^\circ$ being injective (\cref{thm:gluing}(b)).

(b) A lift is $\Theta_\vartheta y$, and $U_\vartheta\Theta_\vartheta y$ is a positive multiple of $\mathrm{glue}_\vartheta(\Theta^\circ y)$, which is orthogonal to the false lifts by \cref{thm:gluing}(d).

(c) Take the frame of \cref{lem:tight-frame}(b) at representatives integral above the prime $\ell$. The forms $\varphi'_a=\mathrm{glue}_\vartheta f_a$ are integral at $\iota$ and at its conjugate (\cref{thm:gluing}(c)), the gauge at a rational translate being the gauge
of the translated vector. For an integral $\varphi$ the coefficients $n_{\rm Aut}\langle\varphi'_a,\varphi\rangle$ are integral, the inner product at a free level being a sum over the classes of pairings of integral vectors and $n_{\rm Aut}$ clearing the
orders of the stabilisers. By \cref{lem:tight-frame}(c), $U_\vartheta\varphi$ is a sum of integral multiples of integral forms.

(d) $U_1\varphi$ is glued from the tuple $\sum_ac_af_a$ with $c_a=n_{\rm Aut}\langle\mathrm{glue}\,f_a,\varphi\rangle$. By Cauchy--Schwarz,
\[
  \abs{\sigma c_a}^2\le n_{\rm Aut}^2\norm{\mathrm{glue}\,f_a}_\sigma^2\norm\varphi_\sigma^2\le n_{\rm Aut}^2\exp(2fmb\log^2b)h_k^2\norm\varphi_\sigma^2,
\]
and the tuple
$\sum c_af_a$ has squared norm at most $(\#\text{frame})^2\max\abs{\sigma c_a}^2h_k^2$; gluing multiplies by at most $\exp(2fmb\log^2b)$. With \cref{lem:tight-frame}(c) the product of these factors is at most $\exp(6fmb\log^2b)$.
\end{proof}

The comparison operator bounds the cost of one lift. Let $\mathrm{pr}$ be the orthogonal projection to the lifts and $\mathrm{pr}\,U_\vartheta$ the compression. Then $\det(\Theta_\vartheta^*\Theta_\vartheta)\det(\mathrm{pr}\,U_\vartheta)=C\rho_\vartheta^d$ with $C\ne0$ independent of $\vartheta$ and $d$
the dimension of the good forms: for an injective $u$ and an operator $T$, $\det(\mathrm{pr}\,T)$ on the range of $u$ times $\det(u^*u)$ is $\det(u^*Tu)$, because $u^*$ kills the orthogonal complement of the range, and $u^*U_\vartheta u=\rho_\vartheta U^\circ$ by (a).
The lifts lie in the image of $e_0$ and are killed by the Hecke polynomial at the root of every split place off the level (\cref{thm:norm-one-lift}(d)), the lift is injective, and for a form orthogonal to the false lifts the projection to the
lifts does not raise its norm at any $\sigma$. A determinant is the product of the eigenvalues; if an operator $\Lambda$ on the lifts with $\abs{\iota(\det\Lambda)}\ge\exp(-dr\beta)$ makes $\Lambda\,\mathrm{pr}\,y$ integral for every integral $y$ orthogonal to the false lifts, then
$\Lambda\,\mathrm{pr}\,U_\vartheta$ keeps the integral lifts by (b), (c), its eigenvalues are integral, and $\abs{\iota\det(\mathrm{pr}\,U_\vartheta)}\le\exp(dr\beta)$; at the trivial twist the conjugates of $\det(\mathrm{pr}\,U_1)$ are at most $\exp(3dfmb\log^2b)$ by (d). These
bounds give the height of $Z_j(1)/Z_j(\xi^j_\iota)$, which is $\det(\mathrm{pr}\,U_\xi)/\det(\mathrm{pr}\,U_1)$ times the $d$-th power of the inverse ratio of Hecke values; this is used in \cref{sec:theorem-C}.
